\chapter{Viewing and Processing Text}\label{ch:sh-text}

Unix's founding bet was that if every program reads and writes plain text, then
a small number of general tools can be combined to solve problems nobody
anticipated. The tools in this chapter are that bet's payoff. Each does one
narrow thing; \autoref{ch:sh-pipes} is where they become a language.

\section{Reading Files}\label{sec:sh-view}

\begin{keys}[0.24]
	\cmd{cat file}        & Print the whole file                                       \\
	\cmd{cat -n file}     & \dots{} with line numbers                                  \\
	\cmd{cat a b > c}     & Concatenate --- which is what the name is short for        \\
	\cmd{head file}       & The first ten lines                                        \\
	\cmd{head -n 3 file}  & The first three                                            \\
	\cmd{tail file}       & The last ten lines                                         \\
	\cmd{tail -n +5 file} & From line five to the end                                  \\
	\cmd{tail -f log}     & Print new lines as they are appended, forever              \\
	\cmd{less file}       & Page through it interactively                              \\
\end{keys}

\cmd{less} is a pager, and its keys will look familiar:

\begin{keys}[0.20]
	\key{<Space>} / \key{b}    & Forward / back one screen                      \\
	\key{j} \key{k}            & Down / up one line                             \\
	\key{g} / \key{G}          & Start / end of file                            \\
	\key{/}\texttt{pat}, \key{n}, \key{N} & Search, next match, previous match  \\
	\key{-N}                   & Toggle line numbers                            \\
	\key{F}                    & Follow, like \cmd{tail -f}; \key{<C-c>} stops  \\
	\key{q}                    & Quit                                           \\
\end{keys}

\begin{note}
	These are Vim's keys, not a coincidence: \cmd{less}, \cmd{man}, \cmd{git log}
	and most things that paginate use the same bindings. Part~\ref{part:vim} has
	already taught you how to read a log file.
\end{note}

\begin{moral}
	\cmd{cat file | grep x} works but wastes a process; \cmd{grep x file} is the
	same thing. This gets called ``useless use of \cmd{cat}'' often enough to be a
	running joke. It is harmless --- but \cmd{cat} on a 2\,GB log is not, and
	\cmd{less} is what you wanted.
\end{moral}

\section{Slicing Lines}\label{sec:sh-slice}

\begin{keys}[0.28]
	\cmd{wc -l file}            & Count lines; \cmd{-w} words, \cmd{-c} bytes       \\
	\cmd{cut -d: -f1 /etc/passwd} & Field 1, with \texttt{:} as the delimiter       \\
	\cmd{cut -c1-10 file}       & Characters 1 to 10 of every line                  \\
	\cmd{paste a b}             & Join two files side by side, line for line        \\
	\cmd{tr 'a-z' 'A-Z'}        & Translate characters                              \\
	\cmd{tr -d '\textbackslash r'} & Delete characters --- this one strips Windows line endings \\
	\cmd{tr -s ' '}             & Squeeze repeats into one                          \\
	\cmd{column -t}             & Align whitespace-separated input into columns     \\
	\cmd{nl file}               & Number the lines                                  \\
	\cmd{rev}                   & Reverse each line's characters                    \\
\end{keys}

\begin{note}
	\cmd{tr} reads only from standard input --- it takes no filename. This is
	deliberate: it is a \vocab{filter}, meant to sit in the middle of a pipeline
	(\autoref{sec:sh-pipes}), and \cmd{tr -d '\textbackslash r' < old > new} is how
	you feed it a file.
\end{note}

\begin{eg}
	The users on this system, one per line:
\begin{shell}
$ cut -d: -f1 /etc/passwd | head -3
root
daemon
bin
\end{shell}
	\cmd{cut} is the right tool exactly when fields are separated by a single
	character. When they are separated by \emph{runs} of whitespace, it is the
	wrong tool and \cmd{awk} is the right one.
\end{eg}

\section{Ordering}\label{sec:sh-sort}

\begin{keys}[0.24]
	\cmd{sort file}          & Sort lines alphabetically                          \\
	\cmd{sort -n}            & Numerically --- \texttt{10} after \texttt{9}, not before \\
	\cmd{sort -h}            & By human-readable size: \texttt{2K}, \texttt{1M}, \texttt{3G} \\
	\cmd{sort -r}            & Reverse                                            \\
	\cmd{sort -u}            & Sort and remove duplicates                         \\
	\cmd{sort -k2}           & By the second field                                \\
	\cmd{sort -t: -k3 -n}    & By field 3, colon-separated, numerically           \\
	\cmd{uniq}               & Collapse \emph{adjacent} duplicate lines           \\
	\cmd{uniq -c}            & \dots{} and prefix each with its count             \\
	\cmd{uniq -d}            & Only the lines that were duplicated                \\
	\cmd{shuf}               & Shuffle; \cmd{shuf -n 1} picks one at random       \\
\end{keys}

\begin{note}
	\cmd{uniq} only removes \emph{adjacent} duplicates, so it is nearly always
	preceded by \cmd{sort}. This looks like a design flaw and is not: it makes
	\cmd{uniq} a streaming filter that needs no memory, and the sort you would
	have needed anyway is written explicitly.
\end{note}

\begin{eg}[The counting idiom]
	This pipeline appears everywhere, and is worth recognising on sight:
\begin{shell}
$ sort access.log | uniq -c | sort -rn | head
   842 /index.html
   311 /about
    97 /contact
\end{shell}
	Sort so duplicates become adjacent; count them; sort those counts numerically,
	largest first; take the top few. \emph{What are the most frequent lines?} ---
	in four words.
\end{eg}

\section{Comparing}\label{sec:sh-compare}

\begin{keys}[0.24]
	\cmd{diff a b}        & Line-by-line differences                              \\
	\cmd{diff -u a b}     & Unified format --- what Git and \cmd{patch} both speak \\
	\cmd{diff -r d1 d2}   & Recursively, over two directories                     \\
	\cmd{cmp a b}         & Byte-for-byte; silent when identical                  \\
	\cmd{comm -12 a b}    & Lines common to two \emph{sorted} files               \\
	\cmd{patch < fix.diff} & Apply a diff                                         \\
	\cmd{vimdiff a b}     & The same comparison, editable (\autoref{sec:vim-diff}) \\
\end{keys}

\begin{eg}[Reading a unified diff]
\begin{shell}
$ diff -u old.txt new.txt
--- old.txt
+++ new.txt
@@ -1,3 +1,3 @@
 unchanged context line
-a line that was removed
+a line that was added
 unchanged context line
\end{shell}
	\texttt{-{}-{}-} is the old file and \texttt{+++} the new; \texttt{@@} gives
	the line ranges of the \vocab{hunk}; then \texttt{-} for removed, \texttt{+}
	for added, and a leading space for context. This is exactly what
	\cmd{git diff} prints, so \autoref{sec:git-diff} needs no further explanation.
\end{eg}

\section{Regular Expressions}\label{sec:sh-regex}

\begin{definition}[Regular expression]
	A pattern language for matching text. Unlike a glob
	(\autoref{sec:sh-glob}), which is expanded by the shell against filenames, a
	regular expression is interpreted by the \emph{program} you hand it to, and
	matches anywhere inside a line.
\end{definition}

The unhappy part is that there are three dialects in common use:

\begin{keys}[0.24]
	\vocab{BRE} & Basic. \cmd{grep} and \cmd{sed} by default. \texttt{+}, \texttt{?}, \texttt{|}, \texttt{(} must be backslash-escaped to have their special meaning \\
	\vocab{ERE} & Extended. \cmd{grep -E}, \cmd{sed -E}, \cmd{awk}. Those characters work unescaped \\
	\vocab{PCRE} & Perl-compatible. \cmd{grep -P}, \cmd{rg}, and most programming languages. Adds \texttt{\textbackslash d}, \texttt{\textbackslash w}, lookaround, non-greedy \texttt{*?} \\
\end{keys}

\begin{moral}
	Use \cmd{grep -E} and \cmd{sed -E} habitually. The escaping rules of BRE exist
	for compatibility with 1975 and cost you a backslash every few characters for
	no benefit.
\end{moral}

The atoms are common to all three:

\begin{keys}[0.20]
	\texttt{.}                & Any single character                           \\
	\texttt{*}                & Zero or more of the preceding item             \\
	\texttt{+}                & One or more                                    \\
	\texttt{?}                & Zero or one                                    \\
	\texttt{\{2,5\}}          & Between two and five                           \\
	\texttt{[abc]}, \texttt{[a-z]} & One character from a set or range         \\
	\texttt{[\textasciicircum{}abc]} & One character not in the set            \\
	\texttt{\textasciicircum{}} & Start of line                                \\
	\texttt{\$}               & End of line                                    \\
	\texttt{(ab|cd)}          & Group, and alternation                         \\
	\texttt{\textbackslash .} & A literal dot                                  \\
	\texttt{\textbackslash b} & A word boundary                                \\
\end{keys}

\begin{note}[The same symbol, two meanings]
	In a glob, \texttt{*} means ``any run of characters''. In a regular expression
	it means ``any number of \emph{the thing before me}''. So the glob
	\texttt{*.txt} is the regular expression
	\texttt{.*\textbackslash .txt\$} --- and \texttt{*.txt} as a regular
	expression means something else entirely. Mixing these up is the single most
	common regular-expression mistake.
\end{note}

\begin{remark}
	Regular expressions are greedy: they match as much as they can.
	\texttt{<.*>} against \texttt{<a> and <b>} matches the whole string, not
	\texttt{<a>}. In PCRE, \texttt{<.*?>} is the non-greedy form; in ERE the usual
	fix is to exclude the delimiter, \texttt{<[\textasciicircum{}>]*>}.
\end{remark}

\section{\texorpdfstring{\cmd{grep}}{grep}}\label{sec:sh-grep}

\cmd{grep} prints the lines of its input that match a pattern. That is all it
does, and it is probably the command you will run most often for the rest of
your career.

\begin{keys}[0.24]
	\cmd{grep pat file}     & Lines containing \texttt{pat}                       \\
	\cmd{grep -i}           & Ignore case                                         \\
	\cmd{grep -v}           & Invert: lines that do \emph{not} match              \\
	\cmd{grep -n}           & Prefix each with its line number                    \\
	\cmd{grep -r pat dir}   & Recurse into a directory                            \\
	\cmd{grep -l}           & Print only the names of matching files              \\
	\cmd{grep -c}           & Print only a count                                  \\
	\cmd{grep -w}           & Match whole words only                              \\
	\cmd{grep -F}           & Fixed strings: no pattern language at all           \\
	\cmd{grep -E}           & Extended regular expressions                        \\
	\cmd{grep -A3 -B1}      & Show 3 lines after and 1 before each match          \\
	\cmd{grep -o}           & Print only the matching part, not the whole line    \\
\end{keys}

\begin{eg}
\begin{shell}
$ grep -rn 'TODO' --include='*.tex' .        # every TODO in the note, with line numbers
$ grep -c '^$' report.tex                    # how many blank lines?
$ grep -o '[a-z]*@[a-z.]*' contacts.txt      # pull out just the addresses
$ ps aux | grep -v grep | grep nvim          # find a process, minus grep itself
\end{shell}
\end{eg}

\begin{note}
	\cmd{grep} sets its exit status to 0 if it matched and 1 if it did not, which
	makes it usable as a test: \cmd{if grep -q pat file; then ...} with \cmd{-q}
	for silence. \autoref{sec:sh-status} builds on this.
\end{note}

\begin{remark}
	\cmd{rg} (ripgrep, installed here) is \cmd{grep -r} done properly: it is
	several times faster, recurses by default, skips anything in
	\file{.gitignore}, and uses PCRE-ish syntax throughout. For searching a source
	tree, prefer it. For a pipeline or a script that must run anywhere, use
	\cmd{grep} --- it is on every Unix machine in existence.
\end{remark}

\section{\texorpdfstring{\cmd{sed}}{sed}}\label{sec:sh-sed}

\cmd{sed} is a \vocab{stream editor}: it reads a line, applies your commands,
prints the result, and repeats. You will use perhaps three of its features ever,
and one of them constantly.

\begin{center}
	\ttfamily\large
	sed [-E] [-n] [-i] '\{address\}\{command\}' \{file\}
\end{center}

\begin{keys}[0.30]
	\cmd{sed 's/old/new/'}       & Replace the first match on each line          \\
	\cmd{sed 's/old/new/g'}      & Replace every match                           \\
	\cmd{sed 's/old/new/2'}      & Replace only the second match on each line    \\
	\cmd{sed -E 's/(a)(b)/\textbackslash 2\textbackslash 1/'} & Swap two capture groups \\
	\cmd{sed -n '5p'}            & Print only line 5 (\cmd{-n} suppresses the default printing) \\
	\cmd{sed -n '10,20p'}        & Print lines 10 to 20                          \\
	\cmd{sed '/pat/d'}           & Delete matching lines                         \\
	\cmd{sed '1d'}               & Delete the first line --- a header, usually   \\
	\cmd{sed -i.bak 's/a/b/g'}   & Edit the file in place, keeping a backup      \\
\end{keys}

\begin{note}
	If this looks familiar, it should: Vim's \cmd{:\%s/old/new/g}
	(\autoref{sec:vim-subst}) and \cmd{:g/pat/d} (\autoref{sec:vim-global}) are the
	same commands with the same syntax. Both inherited them from \texttt{ed}, and
	learning either teaches you the other.
\end{note}

\begin{note}[\cmd{-i} is not portable]
	GNU \cmd{sed} accepts \cmd{sed -i 's/a/b/'}; BSD and macOS \cmd{sed} require an
	argument, \cmd{sed -i '' 's/a/b/'}. Writing \cmd{sed -i.bak} works on both and
	leaves you a backup, which is the right answer for a command that edits files
	irreversibly.
\end{note}

\begin{eg}
	The delimiter need not be a slash --- any character works, which spares you
	from escaping paths:
\begin{shell}
$ sed 's|/usr/local|/opt|g' config.ini
\end{shell}
\end{eg}

\section{\texorpdfstring{\cmd{awk}}{awk}}\label{sec:sh-awk}

\cmd{awk} is a complete programming language, but you can get most of its value
from one sentence: \emph{for each line, split it into fields, and if the line
	matches a condition, do something.}

\begin{center}
	\ttfamily\large
	awk '\{pattern\} \{ \{action\} \}' \{file\}
\end{center}

Either part may be omitted: with no pattern the action runs on every line; with
no action the line is printed.

\begin{keys}[0.20]
	\texttt{\$0}         & The whole line                                        \\
	\texttt{\$1}, \texttt{\$2} & The first and second fields                     \\
	\texttt{NF}          & The number of fields --- so \texttt{\$NF} is the last  \\
	\texttt{NR}          & The current line number                               \\
	\texttt{-F:}         & Use \texttt{:} as the field separator                 \\
	\texttt{BEGIN\{\}}   & Run before the first line                             \\
	\texttt{END\{\}}     & Run after the last                                    \\
\end{keys}

\begin{eg}
\begin{shell}
$ awk '{print $1}' access.log            # first field of every line
$ awk -F: '$3 >= 1000 {print $1}' /etc/passwd    # human users only
$ awk 'NF' file.txt                      # drop blank lines: NF is 0, which is false
$ awk '{s += $2} END {print s}' data     # sum the second column
$ awk 'length > 79' report.tex           # lines over 79 characters
\end{shell}
\end{eg}

\begin{intuition}
	Choose between the three by what the input looks like. If you want
	\emph{lines that match}, use \cmd{grep}. If you want \emph{the same
		substitution everywhere}, use \cmd{sed}. If the input has \emph{columns} ---
	and you want to compare, compute with, or reorder them --- use \cmd{awk}. Most
	one-liners that get uncomfortable in \cmd{sed} are two clear words of
	\cmd{awk}.
\end{intuition}

\begin{moral}
	When an \cmd{awk} program grows past one line, stop and write it in a real
	language. The tools in this chapter are for the middle of a pipeline; they are
	not a place to build software.
\end{moral}
